\documentclass[sigconf,review,anonymous]{acmart}

\microtypesetup{expansion=false}
\settopmatter{printacmref=false}
\renewcommand\footnotetextcopyrightpermission[1]{}
\pagestyle{plain}
\emergencystretch=2em

\title{ArenaWealth Pro: Reproducible Moat and Compounding Analytics for Long-Horizon Equity Portfolios}

\author{Anonymous Author(s)}
\affiliation{%
  \institution{Anonymous Institution}
  \city{Anonymous}
  \country{Anonymous}
}
\email{anonymous@example.com}

\begin{abstract}
Long-horizon equity investors often combine qualitative moat judgments with
quantitative evidence about capital efficiency, free-cash-flow durability, and
valuation. We present ArenaWealth Pro, a reproducible portfolio-analysis
artifact that turns those criteria into an auditable cash-deployment workflow.
The current artifact provides deterministic offline analysis, live provider
adapters, theme-aware concentration controls, and reviewer scripts that validate
the executable surface. This draft is intentionally conservative: it reports
artifact readiness and defines the empirical work still required before making
outperformance claims.
\end{abstract}

\keywords{computational finance, portfolio analytics, reproducibility, moat, compounding}

\begin{document}
\maketitle

\section{Introduction}

ArenaWealth Pro is scoped for ACM-style AI finance venues such as ICAIF
\cite{acmIcaif}. It targets a practical investment workflow: deciding where to
deploy incremental cash inside a concentrated portfolio of high-quality
businesses. The system ranks existing holdings using moat, compounding,
valuation, concentration, and theme constraints, then emits a fee-aware purchase
plan. The artifact is designed to be reviewed as software first: a reviewer can
install dependencies, run tests, reproduce deterministic recommendations, and
inspect claim-evidence mappings.

\section{Scope and Claims}

The current artifact does not claim to beat market benchmarks or registered
advisers. Those claims require out-of-sample backtests, baselines, ablations,
transaction-cost modeling, and licensed data access. The present contribution is
an auditable implementation substrate for those experiments.

\section{System Design}

The backend separates API DTOs, portfolio importers, provider adapters,
analytics scoring, and deployment planning. Network access is isolated behind
provider abstractions. The scoring and deployment logic are deterministic and
covered by unit tests. The browser application is read-only and does not place
orders.

\section{Reviewer Workflow}

The reviewer executes:

\begin{verbatim}
make verify
make metrics
python scripts/moat_compounding_analysis.py \
  --cash 1511.18 \
  --holdings tests/fixtures/seed_portfolio_avenue.csv \
  --offline-demo
\end{verbatim}

The first command validates lint, Python tests, frontend build, and frontend
lint. The second emits a JSON evidence report under \texttt{exports/}. The third
runs a deterministic recommendation path that does not require API keys. The
scaffold follows ACM conference authoring conventions \cite{acmAuthorSubmissions}.

\section{Experimental Plan}

The next paper milestone is a full empirical section:

\begin{itemize}
  \item Baselines: market-cap index, equal weight, quality factor, value factor,
  and simple robo-advisor rebalancing rules.
  \item Ablations: moat, compounding, valuation, concentration, and theme caps.
  \item Costs: fees, slippage, turnover, taxes, and cash drag.
  \item Robustness: provider missingness, stale data, and threshold sensitivity.
\end{itemize}

\section{Limitations}

ArenaWealth Pro is not a registered investment adviser, broker-dealer, or order
execution system. Commercial use requires legal, compliance, data-licensing,
security, suitability, and disclosure review. The repository also does not
include MITRE/STICKS cybersecurity campaign execution; those measurements belong
to a separate artifact.

\section{Conclusion}

ArenaWealth Pro is a cleaned, executable foundation for a computational-finance
artifact. It is ready for reviewer-oriented software validation, while its
research claims remain intentionally bounded until the empirical benchmark suite
is implemented.

\bibliographystyle{ACM-Reference-Format}
\bibliography{references}

\end{document}
